\begin{theorem}[Uniform interior $C^{k,\beta}$ estimates for planar $p$-harmonic functions]
Let $1<p<2$, let the Lewis rigidity datum of \noderef{LewisRigidityInput}, the Lewis--Tolksdorf
datum $F$ of \noderef{FirstOrderEstimateInput}, the Iwaniec--Manfredi datum of
\noderef{IMRegularityInput}, the smoothness-away-from-critical-points hypothesis of
\noderef{SmoothAwayFromCriticalInput}, the interior divergence-form Schauder estimates of
\noderef{DivFormSchauderInput} and the local boundedness datum of
\noderef{LocalBoundednessInput} be given for $p$.  Let $k\in\mathbb{N}$ with $k\ge2$ and let
$0<\beta<1$ satisfy $k+\beta<s(p)$, with $s(p)$ as in \noderef{SOpt}.

Then there is $C>0$, depending only on $p$, $k$, $\beta$ and the data, with the following
property.  Let $u\colon\mathrm{Plane}\to\mathbb{R}$ be $p$-harmonic on $B_1$ with weak gradient
$g$ in the sense of \noderef{IsPHarm} and suppose $u\in L^{p}(B_1)$, i.e.\ $u$ is
$p$-integrable with respect to Lebesgue measure restricted to $B_1$; write
$N:=\lVert u\rVert_{L^{p}(B_1)}=\bigl(\int_{B_1}\lvert u\rvert^{p}\,dx\bigr)^{1/p}$.  Then:
\begin{enumerate}
\item for every $j$ with $0\le j\le k$, every $x\in B_{1/2}$ and every $w\in\mathrm{Plane}$,
\[
  \bigl\lvert D^{j}u(x)[w,\ldots,w]\bigr\rvert\le C\,N\,\lvert w\rvert^{j};
\]
\item for all $x,y\in B_{1/2}$ and every $w\in\mathrm{Plane}$,
\[
  \bigl\lvert D^{k}u(x)[w,\ldots,w]-D^{k}u(y)[w,\ldots,w]\bigr\rvert
  \le C\,N\,\lvert x-y\rvert^{\beta}\,\lvert w\rvert^{k}.
\]
\end{enumerate}
Here $D^{j}u(x)$ is the $j$-th iterated Fr\'echet derivative of $u$ at $x$, a continuous
$j$-linear form on $\mathrm{Plane}$, evaluated on $j$ copies of $w$; at $j=0$ it is the
constant $0$-linear form with value $u(x)$, so clause (1) at $j=0$ reads
$\lvert u(x)\rvert\le C\,N$.  That $C$ is produced before $u$ is the whole content of the
theorem: it is the uniformity the paper's abstract asserts and that the qualitative higher
regularity of Iwaniec--Manfredi does not supply.

\medskip\noindent\emph{How this statement renders the paper's (1.3), and why no deviation is
claimed.}  The paper's (1.3) (line 111) reads
\[
  \sum_{j=0}^{k}\bigl\lVert D^{j}u\bigr\rVert_{L^{\infty}(B_{1/2})}
    +\bigl[D^{k}u\bigr]_{C^{0,\beta}(B_{1/2})}\le C\lVert u\rVert_{L^{p}(B_1)} .
\]
Two things separate that from the two clauses above, and both are reconciled by a factor
depending only on $k$.

\emph{(a) Diagonal evaluations versus full norms.}  (1.3) is written with the operator norms
$\lVert D^{j}u(x)\rVert$ and with the seminorm of the full form $D^{k}u$, whereas clauses (1)
and (2) carry \emph{diagonal} evaluations $D^{j}u(x)[w,\ldots,w]$, homogeneous of degree $j$ in
$w$.  The two directions are not symmetric in difficulty.  One direction is the background
proposition on the operator norm of a continuous multilinear map, which gives
$\lvert Q[v_1,\ldots,v_m]\rvert\le\lVert Q\rVert\lvert v_1\rvert\cdots\lvert v_m\rvert$ and
hence, on the diagonal with $\lvert w\rvert\le1$, that a full-norm bound implies the diagonal
one.  The other direction is the one (1.3) needs: for a \emph{symmetric} $m$-linear form $T$
one has $\lVert T\rVert\le(m^{m}/m!)\sup_{\lvert w\rvert\le1}\lvert T[w,\ldots,w]\rvert$, so
taking the supremum over $\lvert w\rvert\le1$ in clauses (1) and (2) recovers
$\lVert D^{j}u\rVert_{L^{\infty}(B_{1/2})}$ and $[D^{k}u]_{C^{0,\beta}(B_{1/2})}$, bounded by
$C\,N$ times a factor $m^{m}/m!$ depending only on $m\le k$.  The symmetry hypothesis is not a
gap in the mathematics: for $u$ of class $C^{j}$ the form $D^{j}u(x)$ \emph{is} symmetric, by
the classical higher-order Schwarz--Clairaut theorem, and a symmetric form is determined by its
diagonal.  What is unavailable is Mathlib's general-order symmetry of the iterated derivative
for finitely smooth functions, which at the pinned revision exists only at order two --- the
general-order permutation invariance is available there only under analyticity, which none of the
data above supplies --- and the
counterexample recorded against the diagonal-to-full direction, the trilinear
$A[x,y,z]=x_1y_1z_2-x_1y_2z_1$ with $A[v,v,v]=0$ for every $v$ but $A[e_1,e_1,e_2]=1$, is
\emph{not} symmetric, so it shows only that the recovery genuinely needs symmetry, not that it
fails for the forms at hand.  The diagonal rendering is therefore a reachability choice, not a
weakening.

\emph{(b) Two families versus one sum.}  (1.3) is a single inequality on a sum of $k+2$
terms --- the $k+1$ supremum norms indexed by $j=0,\ldots,k$ inclusive, plus the one H\"older
seminorm --- whereas clauses (1) and (2) are two quantified families of bounds, one bound per
term.  Each direction is immediate.  Forwards: if each of the $k+2$ terms is at most $C\,N$,
then their sum is at most $(k+2)\,C\,N$, so (1.3) holds with the constant $(k+2)\,C$.
Backwards: every term of the sum is non-negative, so the sum dominates each term individually,
and (1.3) with constant $C$ gives each clause with the same $C$.  The family rendering is
thus neither weaker nor stronger than the paper's sum, up to the factor $k+2$.

Both reconciling factors, $m^{m}/m!$ for $m\le k$ in (a) and $k+2$ in (b), depend only on $k$,
and the paper's own constant is quantified as $C=C(p,k,\beta)$ --- line 313 states outright that
all constants depend only on $p,k,\beta$ --- so both are absorbed by it and the uniformity that
is the content of (1.3) is untouched.  That is why these two renderings are recorded here as
part of the statement rather than claimed as a deviation: the paper's statement is
mathematically correct, and the difference is a repackaging together with one
Mathlib-reachability artefact.
\end{theorem}
\begin{proof}
\medskip\noindent\textbf{Step 1: the constant, fixed before $u$.}
Apply \noderef{JetBounds} with $p$, $1<p<2$, the six data, $k$, $k\ge2$, $\beta$, $0<\beta<1$
and $k+\beta<s(p)$: it produces $C_1>0$ such that every $p$-harmonic $u$ on $B_1$ with weak
gradient $g$ and $u\in L^{p}(B_1)$ satisfies, for every $x\in B_{1/2}$,
\begin{equation}\label{eq:main-jb}
  \bigl\lvert D^{j}u(x)[w,\ldots,w]\bigr\rvert\le C_1N\lvert w\rvert^{j}
  \qquad(1\le j\le k,\ w\in\mathrm{Plane}),
\end{equation}
together with an error bound we do not use here.  Apply \noderef{HolderFromJets} with the same
arguments: it produces $C_2>0$ such that every such $u$ satisfies
\begin{equation}\label{eq:main-hfj}
  \bigl\lvert D^{k}u(x)[w,\ldots,w]-D^{k}u(y)[w,\ldots,w]\bigr\rvert
  \le C_2N\lvert x-y\rvert^{\beta}\lvert w\rvert^{k}
  \qquad(x,y\in B_{1/2},\ w\in\mathrm{Plane}).
\end{equation}
Let $C_b>0$ be the constant carried by the local boundedness datum of
\noderef{LocalBoundednessInput}.  Put
\[
  C:=\max\bigl\{C_1,\ \max\{C_2,\ C_b\}\bigr\}.
\]
Then $C>0$, since $C_1>0$ and $C_1\le C$; and $C_1\le C$, $C_2\le C$ and $C_b\le C$.  All three
of $C_1,C_2,C_b$ were produced from $p,k,\beta$ and the data alone, before any $u$, so $C$ was
too.

Now fix $u$ and $g$ as in the statement and put $N:=\lVert u\rVert_{L^{p}(B_1)}$.  Note
$N\ge0$, being an $L^{p}$ norm.  Note also $B_{1/2}\subseteq B_{3/4}$, since
$\lvert x\rvert<1/2$ implies $\lvert x\rvert<3/4$.

\medskip\noindent\textbf{Step 2: clause (1) at $j=0$, which is local boundedness.}
Let $x\in B_{1/2}$ and $w\in\mathrm{Plane}$.  By clause (1) of the background proposition on
the recursion for the iterated Fr\'echet derivative, $D^{0}u(x)$ is the constant $0$-linear
form with value $u(x)$, so its evaluation on the empty family of vectors is $u(x)$ and
\[
  \bigl\lvert D^{0}u(x)[\,]\bigr\rvert=\lvert u(x)\rvert .
\]
Moreover $\lvert w\rvert^{0}=1$, this being the empty product.  So the assertion of clause (1)
at $j=0$ is exactly $\lvert u(x)\rvert\le C\,N$.

Since $x\in B_{1/2}\subseteq B_{3/4}$, the clause \texttt{bound} of
\noderef{LocalBoundednessInput}, applied to $u$, $g$, the $p$-harmonicity hypothesis and
$u\in L^{p}(B_1)$, gives $\lvert u(x)\rvert\le C_b\,N$.  As $C_b\le C$ and $N\ge0$,
multiplying the inequality $C_b\le C$ by $N$ preserves it, so $C_b\,N\le C\,N$ and therefore
$\lvert u(x)\rvert\le C\,N$.  This is clause (1) at $j=0$.

\medskip\noindent\textbf{Step 3: clause (1) for $1\le j\le k$, and clause (2).}
Let $1\le j\le k$, let $x\in B_{1/2}$ and let $w\in\mathrm{Plane}$.  By \eqref{eq:main-jb},
\[
  \bigl\lvert D^{j}u(x)[w,\ldots,w]\bigr\rvert\le C_1N\lvert w\rvert^{j}.
\]
Now $C_1\le C$ and $N\ge0$ give $C_1N\le C\,N$, and $\lvert w\rvert^{j}\ge0$, being a power of a
non-negative real, so multiplying through gives $C_1N\lvert w\rvert^{j}\le C\,N\lvert
w\rvert^{j}$.  Hence $\lvert D^{j}u(x)[w,\ldots,w]\rvert\le C\,N\lvert w\rvert^{j}$.  Together
with Step 2, which covers $j=0$, clause (1) holds for every $j\le k$.

For clause (2), let $x,y\in B_{1/2}$ and $w\in\mathrm{Plane}$.  By \eqref{eq:main-hfj},
\[
  \bigl\lvert D^{k}u(x)[w,\ldots,w]-D^{k}u(y)[w,\ldots,w]\bigr\rvert
  \le C_2N\lvert x-y\rvert^{\beta}\lvert w\rvert^{k}.
\]
Here $\lvert x-y\rvert^{\beta}\ge0$: the base is non-negative, and by clause (3) of the
background proposition on elementary properties of real powers the map $t\mapsto t^{\beta}$ is
nondecreasing on $[0,\infty)$, while clause (1) of that proposition gives $0^{\beta}=0$ because
$\beta\ne0$; so $\lvert x-y\rvert^{\beta}\ge0^{\beta}=0$.  Likewise
$\lvert w\rvert^{k}\ge0$.  From $C_2\le C$ and $N\ge0$ we get $C_2N\le C\,N$, and multiplying
successively by the non-negative factors $\lvert x-y\rvert^{\beta}$ and $\lvert w\rvert^{k}$
preserves the inequality, so
$C_2N\lvert x-y\rvert^{\beta}\lvert w\rvert^{k}\le C\,N\lvert x-y\rvert^{\beta}\lvert
w\rvert^{k}$.  Hence clause (2) holds.

\medskip\noindent As $u$ and $g$ were arbitrary subject to the stated hypotheses, and $C$ was
fixed in Step 1 before them, the proof is complete.
\end{proof}
